\documentclass[a4paper, serif, 10pt]{moderncv}

%% moderncv
% style and color
\moderncvstyle{banking}
\moderncvcolor{black}

%% page layout/margins
\usepackage[top=2.5cm, bottom=2.5cm, left=1.5cm, right=1.5cm]{geometry}

\nopagenumbers{}

%% Babel config for Indonesian (Bahasa Indonesia) language
% ref:
% - https://latex3.github.io/babel/guides/locale-indonesian.html
\usepackage[indonesian]{babel}

%% fontspec
\usepackage{fontspec}

\IfFontExistsTF{Linux Libertine}{
  \setmainfont[Ligatures={TeX, Common}, Numbers=OldStyle]{Linux Libertine}
}{}
\IfFontExistsTF{Linux Libertine O}{
  \setmainfont[Ligatures={TeX, Common}, Numbers=OldStyle]{Linux Libertine O}
}{}

%% CTeX
% CJK support, used to show CJK characters in the resume
%
% - fontset=none: disable builtin fontset but instead set the CJK font manually
% - heading=false: disable ctex heading
% - punct=kaiming: use kaiming punctuations styles for CJK
% - scheme=plain: use plain scheme, do not override \normalsize font size
% - space=auto: space settings for CJK characters
%
% ref:
% - http://ctan.mirrorcatalogs.com/language/chinese/ctex/ctex.pdf
\usepackage[UTF8, heading=false, punct=kaiming, scheme=plain, space=auto]{ctex}

\IfFontExistsTF{Noto Serif CJK SC}{
  \setCJKmainfont{Noto Serif CJK SC}
}{}
\IfFontExistsTF{Noto Sans CJK SC}{
  \setCJKsansfont{Noto Sans CJK SC}
}{}

%% line spacing
\usepackage{setspace}
\setstretch{1.125}

%% URL styling - use normal text instead of monospace
\urlstyle{same}

% Auto-underline all links
% Defer until \begin{document} so hyperref is already loaded
\AtBeginDocument{%
  \let\oldhref\href
  \renewcommand{\href}[2]{\oldhref{#1}{\underline{#2}}}%
}

%% Patch \httplink and \httpslink to handle full URLs
% The original moderncv macros always prepend http:// or https:// to URLs.
% This causes issues when the URL already has a protocol (e.g., https://example.com)
% resulting in malformed URLs like https://https://example.com.
% This patch checks if the URL already contains "://" and uses it directly if so.
% Note: We use \str_set:Nx (x-expansion) to expand macros like \@homepage first.
\ExplSyntaxOn
\str_new:N \g_yamlresume_url_str

\RenewDocumentCommand{\httpslink}{O{}m}{
  \str_set:Nx \g_yamlresume_url_str {#2}
  \str_if_in:NnTF \g_yamlresume_url_str {://}
    { \url{#2} }
    { \url{https://#2} }
}

\RenewDocumentCommand{\httplink}{O{}m}{
  \str_set:Nx \g_yamlresume_url_str {#2}
  \str_if_in:NnTF \g_yamlresume_url_str {://}
    { \url{#2} }
    { \url{http://#2} }
}
\ExplSyntaxOff

\name{Rina Kartika}{}
\title{Product Manager dengan 6+ tahun pengalaman di produk digital \& fintech}
\phone[mobile]{+62 812-3456-7890}
\email{rina.kartika@email.com}
\homepage{https://www.linkedin.com/in/rinakartika}

\address{Jakarta Selatan, DKI Jakarta, Indonesia, 12950}{}{}

\extrainfo{{\small \faLinkedin}\ \href{https://www.linkedin.com/in/rinakartika}{@rinakartika} {} {} {} • {} {} {}
{\small \faGithub}\ \href{https://github.com/rinakartika}{@rinakartika} {} {} {} • {} {} {}
{\small \faMedium}\ \href{https://medium.com/@rinakartika}{@rinakartika}}

\begin{document}

\maketitle

\section{Informasi Dasar}

\cvline{}{\begin{itemize}
\item Product Manager berpengalaman dalam membangun dan mengembangkan produk digital terutama di bidang fintech
\item Terbiasa memimpin siklus hidup produk dari riset pengguna, perencanaan roadmap, hingga peluncuran dan iterasi
\item Memiliki kemampuan kuat dalam analisis data, pengambilan keputusan berbasis metrik, dan kolaborasi lintas fungsi
\item Berpengalaman bekerja dengan tim engineering, desain, pemasaran, dan pemangku kepentingan bisnis
\item Fokus pada peningkatan pengalaman pengguna, pertumbuhan produk, dan pencapaian tujuan bisnis
\end{itemize}}
\section{Pendidikan}

\cventry{Agu 2014–Jul 2018}
        {Sarjana, Manajemen, Nilai: 3.78}
        {Universitas Indonesia}
        {\url{https://www.ui.ac.id/}}
        {}
        {\begin{itemize}
\item Fokus pada strategi bisnis dan manajemen produk digital
\item Aktif dalam organisasi kemahasiswaan dan berbagai studi kasus bisnis
\item Mengembangkan kemampuan analitis dan kepemimpinan melalui proyek kelompok
\end{itemize}
\textbf{Mata Kuliah}: Manajemen Pemasaran, Manajemen Operasi, Sistem Informasi Manajemen, Perilaku Konsumen, Statistika Bisnis, Manajemen Keuangan}
\section{Pengalaman Kerja}

\cventry{Mar 2021–Sekarang}
        {Senior Product Manager}
        {PT Nusantara Digital Solusi}
        {\url{https://www.nusantaradigital.co.id}}
        {}
        {\begin{itemize}
\item Memimpin pengembangan produk fintech dengan lebih dari 2 juta pengguna aktif
\item Menyusun dan mengelola product roadmap serta prioritas fitur berdasarkan data dan kebutuhan pasar
\item Berkolaborasi dengan tim engineering, desain, riset, dan bisnis untuk menghadirkan solusi yang tepat
\item Meningkatkan retensi pengguna sebesar 35\% melalui optimalisasi alur onboarding dan personalisasi layanan
\item Melakukan riset pengguna, A/B testing, dan analisis metrik produk secara berkala
\item Menjadi penghubung antara kebutuhan bisnis dan kapabilitas teknis dalam pengembangan fitur
\end{itemize}
\textbf{Kata Kunci}: Product roadmap, Fintech, Data-driven, Agile, Stakeholder management}

\cventry{Feb 2019–Des 2020}
        {Product Manager}
        {PT Fintek Sejahtera}
        {}
        {}
        {\begin{itemize}
\item Mengelola siklus hidup produk pembayaran digital dari konsep hingga peluncuran
\item Berkoordinasi dengan tim teknik dalam implementasi fitur dan penyelesaian masalah teknis
\item Menganalisis perilaku pengguna untuk mengidentifikasi peluang peningkatan produk
\item Menyusun laporan performa produk dan mempresentasikannya kepada pemangku kepentingan
\item Bekerja sama dengan tim pemasaran dalam menyusun strategi go-to-market
\end{itemize}
\textbf{Kata Kunci}: Pembayaran digital, Product lifecycle, Go-to-market, Analisis data}

\cventry{Jul 2018–Jan 2019}
        {Associate Product Manager}
        {PT Ecommerce Cemerlang}
        {}
        {}
        {\begin{itemize}
\item Mendukung product manager dalam menyusun spesifikasi produk dan user stories
\item Mengelola backlog dan memastikan kelancaran proses sprint bersama tim pengembang
\item Melakukan pengujian produk serta mengumpulkan umpan balik dari pengguna internal
\end{itemize}
\textbf{Kata Kunci}: User stories, Sprint, Backlog, UAT}
\section{Bahasa}

\cvline{Indonesia}{Bahasa Ibu / Bilingual \hfill \textbf{Kata Kunci}: Bahasa Indonesia}
\cvline{Inggris}{Tingkat Profesional Penuh \hfill \textbf{Kata Kunci}: TOEFL iBT 105, IELTS 7.5}
\section{Keahlian}

\cvline{Manajemen Produk}{Ahli \hfill \textbf{Kata Kunci}: Product roadmap, Prioritisasi fitur, Product discovery, Go-to-market strategy, Product lifecycle}
\cvline{Analisis Data}{Mahir \hfill \textbf{Kata Kunci}: Metrik produk, A/B testing, Google Analytics, SQL, Excel}
\cvline{Riset Pengguna}{Mahir \hfill \textbf{Kata Kunci}: Wawancara pengguna, Survei, Persona, Usability testing}
\cvline{Alat dan Metodologi}{Ahli \hfill \textbf{Kata Kunci}: Agile, Scrum, Jira, Figma, Mixpanel}
\section{Penghargaan}

\cventry{Nov 2022}
        {Asosiasi Fintech Indonesia}
        {Penghargaan Inovasi Produk Digital}
        {}
        {}
        {Produk yang dipimpin berhasil meraih penghargaan atas inovasi dalam layanan keuangan digital.}
\section{Sertifikat}

\cventry{Jun 2020}
        {Scrum Alliance}
        {Certified Scrum Product Owner (CSPO)}
        {\url{https://www.scrumalliance.org/}}
        {}
        {}

\cventry{Jul 2022}
        {Product School}
        {Product Management Professional}
        {\url{https://productschool.com/}}
        {}
        {}
\section{Publikasi}

\cventry{Mar 2023}
        {Meningkatkan Adopsi Produk Fintech melalui Pendekatan Berbasis Data}
        {Medium}
        {\url{https://medium.com/@rinakartika/meningkatkan-adopsi-produk-fintech}}
        {}
        {Artikel tentang bagaimana product manager dapat memanfaatkan data perilaku pengguna untuk mendorong pertumbuhan produk digital.}
\section{Proyek}

\cventry{Jan 2022–Des 2022}
        {Aplikasi mobile untuk pengelolaan keuangan pribadi dengan fitur pembuatan anggaran dan pelacakan pengeluaran.}
        {Aplikasi Keuangan DuitKita}
        {\url{https://duitkita.id}}
        {}
        {\begin{itemize}
\item Memimpin pengembangan produk dari riset pengguna hingga peluncuran
\item Berhasil mencapai 100 ribu unduhan dalam enam bulan pertama
\item Menyusun strategi monetisasi dan kolaborasi dengan mitra keuangan
\end{itemize}
\textbf{Kata Kunci}: Fintech, Mobile app, Product strategy}

\cventry{Feb 2020–Okt 2020}
        {Platform loyalitas yang menghubungkan merchant dan pelanggan melalui program poin terintegrasi.}
        {Platform Loyalitas Pelanggan PoinPintar}
        {\url{https://poinpintar.id}}
        {}
        {\begin{itemize}
\item Mengelola kebutuhan produk dan koordinasi lintas tim
\item Menghadirkan fitur redeem poin yang meningkatkan transaksi merchant sebesar 20\%
\end{itemize}
\textbf{Kata Kunci}: Loyalty, Platform, Merchant}
\section{Minat}

\cvline{Teknologi}{Kecerdasan buatan, Produk digital, Startup}
\cvline{Literasi}{Membaca buku, Menulis blog, Podcast produk}
\section{Kegiatan Sukarela}

\cventry{Jan 2023–Sekarang}
        {Mentor Produk}
        {Yayasan Penggerak Ekonomi Digital}
        {\url{https://yayasanpenggerak.id}}
        {}
        {\begin{itemize}
\item Mendampingi startup tahap awal dalam menyusun strategi produk dan validasi pasar
\item Menyampaikan workshop tentang product discovery dan pengukuran metrik produk
\end{itemize}}

\end{document}